\documentclass[12pt]{article}
\usepackage[utf8]{inputenc}
\usepackage[T1]{fontenc}
\usepackage{hyperref}
\title{Sena, Malawi Language Resources for K-12 Curriculum Development}
\author{Compiled by Claude Sonnet 5.5 (claude-sonnet-5-5) for Andr\'as Kornai}
\date{}
\begin{document}
\maketitle

\section{Introduction}
Malawi Sena (ISO 639-3 swk, Glottocode mala1475) is an Atlantic-Congo language of southern Malawi. Joshua Project places the Sena people in the Chikwawa and Nsanje districts. It is coded separately from Mozambican Sena (seh). The searches for this compilation found no school materials in Malawi Sena: no primer, textbook, dictionary or grammar. The list below is therefore short and consists mostly of religious and audio material and general reference pages.

\section{Primary Educational Resources}
No primers, textbooks or school programmes for Malawi Sena turned up in the searches. Anyone building a curriculum will have to create these materials, ideally with community and Malawian education authorities.

\section{Linguistic and Reference Materials}
\begin{itemize}
\item Wikipedia article on the Sena language, which discusses the relation between Malawian Sena and Mozambican Sena: \url{https://wikipedia.com/wiki/Cisena}
\end{itemize}
No dictionary or grammar specific to Malawi Sena was found in the search results.

\section{Religious and Community Resources}
\begin{itemize}
\item Joshua Project language page for swk, which lists scripture and media availability: \url{https://joshuaproject.net/index.php/languages/swk}
\item Joshua Project profile of the Sena, Malawi people group: \url{https://joshuaproject.net/people_groups/14849}
\end{itemize}
The search summaries did not make clear whether a complete Bible exists in this language.

\section{Digital and Audio Resources}
\begin{itemize}
\item Global Recordings Network, audio and audio-visual Bible-teaching materials in Sena, Malawi: \url{https://globalrecordings.net/en/language/swk}
\end{itemize}

\section{Curriculum Development Context}
The language is thought to be spoken by the whole ethnic community, but there is no evidence that it is taught in schools. Existing resources are mostly evangelistic audio recordings, which can serve as a starting point for oral-language material. A mother-tongue curriculum would need new orthography decisions, primers, readers and teacher training. Verify the current orthography status before starting.

\end{document}